\begin{figure}[H]
\centering
\resizebox{\textwidth}{!}{%
\begin{tikzpicture}[
    node distance=1.2cm and 2.0cm,
    >=Stealth, thick, font=\sffamily
]

% ---------------------------------------------------------
% Styles
% ---------------------------------------------------------
\tikzset{
    op_node/.style={
        circle, draw=black, very thick, minimum size=0.9cm, fill=white,
        font=\sffamily\large
    },
    io_node/.style={font=\sffamily\bfseries},
    dot/.style={circle, fill=black, inner sep=0pt, minimum size=4pt}
}

% ---------------------------------------------------------
% CORDIC block (large)
% ---------------------------------------------------------
\node[block, minimum width=4.0cm, minimum height=3.5cm, text width=3.5cm] (CORDIC) at (0, 0) {%
    \textbf{CORDIC}\\[2pt]
    \textsf{\footnotesize cordic\_rect\_to\_polar}\\[4pt]
    {\footnotesize Rect $\rightarrow$ Polar}\\
    {\footnotesize Latency: 40 cyc}
};

% ---------------------------------------------------------
% Inputs
% ---------------------------------------------------------
\node[io_node, left=2.5cm of CORDIC] (Iin) at (-5, 1.0) {$i_{in}$};
\node[io_node] (Qin) at (-5, -1.0) {$q_{in}$};
\node[io_node] (Vin) at (-5, -3.0) {input\_valid};

\draw[data arrow] (Iin) -- ($(CORDIC.west) + (0, 0.6)$)
    node[signal, midway, above] {\scriptsize x\_in};
\draw[data arrow] (Qin) -- ($(CORDIC.west) + (0, -0.6)$)
    node[signal, midway, above] {\scriptsize y\_in};

% ---------------------------------------------------------
% Magnitude path (top)
% ---------------------------------------------------------
% CORDIC x_out -> multiply by K^{-1}
\node[op_node, right=2.0cm of CORDIC] (MultK) at ($(CORDIC.east) + (2.0, 1.0)$) {$\times$};
\draw[data arrow] ($(CORDIC.east) + (0, 1.0)$) -- (MultK)
    node[signal, midway, above] {\scriptsize mag};

% K^{-1} constant from top
\node[const, above=0.8cm of MultK] (Kinv) {$K^{-1} = 0.8588$};
\draw[data arrow] (Kinv) -- (MultK);

% Truncate
\node[small block, right=1.5cm of MultK] (TruncE) {$[47\!:\!16]$};
\draw[data arrow] (MultK) -- (TruncE);

% envelope_out
\node[io_node, right=1.5cm of TruncE] (EnvOut) {envelope\_out};
\draw[data arrow] (TruncE) -- (EnvOut);

% ---------------------------------------------------------
% Phase path (bottom)
% ---------------------------------------------------------
% CORDIC phase_out -> shift right 13
\node[small block] (ShiftR) at ($(CORDIC.east) + (2.0, -1.0)$) {$\gg 13$};
\draw[data arrow] ($(CORDIC.east) + (0, -1.0)$) -- (ShiftR)
    node[signal, midway, above] {\scriptsize phase\_out};

% phase_raw fork point
\coordinate (PhaseRawFork) at ($(ShiftR.east) + (1.0, 0)$);
\node[dot] at (PhaseRawFork) {};
\draw[data arrow] (ShiftR) -- (PhaseRawFork);

% Direct path to phase_wrapped_out
\node[io_node] (PhaseWOut) at ($(PhaseRawFork) + (0, -2.0)$) {phase\_wrapped\_out};
\draw[data arrow] (PhaseRawFork) -- (PhaseWOut);

% Path to Phase Unwrap block
\node[block, minimum width=3.5cm, minimum height=2.5cm, text width=3.0cm] (Unwrap) at ($(PhaseRawFork) + (3.5, 0)$) {%
    \textbf{Phase Unwrap}\\[2pt]
    {\footnotesize $\Delta\theta = \theta(n) - \theta(n\!-\!1)$}\\
    {\footnotesize wrap to $[-\pi, +\pi]$}\\
    {\footnotesize accumulate}
};
\draw[data arrow] (PhaseRawFork) -- (Unwrap.west)
    node[signal, midway, above] {\scriptsize phase\_raw};

% Outputs from Phase Unwrap
\node[io_node, right=1.5cm of Unwrap] (ThetaU) at ($(Unwrap.east) + (2.0, 0.5)$) {theta\_unwrapped\_out};
\node[io_node] (ThetaN) at ($(Unwrap.east) + (2.0, -0.5)$) {theta\_norm\_out};

\draw[data arrow] ($(Unwrap.east) + (0, 0.5)$) -- (ThetaU);
\draw[data arrow] ($(Unwrap.east) + (0, -0.5)$) -- (ThetaN);

% ---------------------------------------------------------
% Valid delay line
% ---------------------------------------------------------
\node[small block, minimum width=3.0cm] (ValDelay) at (0, -4.5) {%
    Valid Delay\\
    {\scriptsize 40-cycle shift register}
};
\draw[data arrow] (Vin) -- ++(1.5, 0) |- (ValDelay.west);

\node[io_node, right=2.5cm of ValDelay] (BufEn) {buffers\_enable};
\draw[data arrow] (ValDelay) -- (BufEn);

% ---------------------------------------------------------
% Entity border
% ---------------------------------------------------------
\node[entity border, fit=(Iin)(Vin)(EnvOut)(ThetaU)(ThetaN)(PhaseWOut)(ValDelay)(BufEn)(Kinv), inner sep=15pt] {};

\end{tikzpicture}%
}

\vspace{3mm}
\begin{align*}
    \text{envelope} &= (\text{mag}_{\text{CORDIC}} \times K^{-1})\,[47\!:\!16] \\
    \theta_{\text{norm}}(n) &= \theta_{\text{raw}}(n) - \theta_{\text{raw}}(n\!-\!1), \quad \text{wrapped to } [-\pi, +\pi] \\
    \theta_{\text{unwrapped}}(n) &= \theta_{\text{unwrapped}}(n\!-\!1) + \theta_{\text{norm}}(n)
\end{align*}

\caption{Input\_Processing\_Block -- CORDIC-based rectangular-to-polar conversion with phase unwrapping.}
\label{fig:input_processing}
\end{figure}
